% Version: 2026.1.0.0
% ============================================================================
% Engineering Studio AI — Multi-Modal Multi-Agent Discipline Framework
% Formal Model and Mathematical Proofs
% Author: Hadrian Hu (hadrian.hu@gmail.com)
% Governed by: coding_stds/documentation/latex_style_guide.txt
%              coding_stds/mathematics/mathematical_proof_structures_standards.txt
% Date: \today
% ============================================================================
\documentclass[12pt]{article}

% ============================================================================
% PREAMBLE
% ============================================================================

% --- Font & Encoding ---------------------------------------------------
\usepackage[utf8]{inputenc}   % WHAT: UTF-8 source encoding (Std 14.1)
\usepackage[T1]{fontenc}      % WHAT: T1 font encoding (Std 14.2)
\usepackage{times}            % WHAT: Times New Roman body font (Std 2.1)

% --- Page Layout ---------------------------------------------------------
\usepackage[margin=1in]{geometry}  % WHAT: 1in margins (Std 2.3)
\usepackage{setspace}
\singlespacing                     % WHY: single spacing throughout (Std 2.4)

% --- Headers & Footers ----------------------------------------------------
\usepackage{fancyhdr}
\usepackage{lastpage}
\pagestyle{fancy}
\fancyhf{}
\rhead{\thepage\ of \pageref{LastPage}}
\lhead{Engineering Studio AI — Formal Model \& Proofs}
\renewcommand{\headrulewidth}{0.4pt}

% --- Math ------------------------------------------------------------------
\usepackage{amsmath}    % WHAT: aligned/multline/equation environments
\usepackage{amssymb}    % WHAT: blacksquare, mathbb, mathcal
\usepackage{amsthm}     % WHAT: theorem/lemma/proof environments (Std proof §4.1.2)
\usepackage{mathtools}  % WHAT: enhanced math tooling, coloneqq

% --- Graphics & Tables -------------------------------------------------------
\usepackage{graphicx}
\usepackage{booktabs}
\usepackage{array}
\usepackage{multirow}
\usepackage{makecell}
\usepackage{float}

% --- Code & Pseudocode -------------------------------------------------------
\usepackage{listings}
\usepackage{algorithm}
\usepackage{algpseudocode}

% --- Lists -------------------------------------------------------------------
\usepackage{enumitem}
\setlist[enumerate]{itemsep=0pt, parsep=0pt}

% --- Typography safety -------------------------------------------------------
\usepackage[final]{microtype}
\setlength{\emergencystretch}{15pt}
\widowpenalty=10000
\clubpenalty=10000

% --- TOC ------------------------------------------------------------------
\usepackage{tocloft}
\setcounter{tocdepth}{2}

% --- Hyperlinks (loaded second-to-last, per Std 27.9) -----------------------
\usepackage[hidelinks]{hyperref}
\hypersetup{
    colorlinks=false,
    pdftitle={Engineering Studio AI: Formal Model and Mathematical Proofs},
    pdfauthor={Hadrian Hu},
    pdfsubject={Multi-Agent Engineering Workflow Formalization}
}

% --- Theorem environments (Std proof §4.2.1) --------------------------------
\theoremstyle{definition}
\newtheorem{theorem}{Theorem}
\newtheorem{lemma}[theorem]{Lemma}
\newtheorem{corollary}[theorem]{Corollary}
\newtheorem{definition}[theorem]{Definition}
\newtheorem{proposition}[theorem]{Proposition}
\newtheorem{remark}[theorem]{Remark}
\newtheorem{hypothesis}{Hypothesis}

% --- Custom commands ---------------------------------------------------------
\newcommand{\code}[1]{\texttt{#1}}
\newcommand{\FN}[1]{\textsc{fn-#1}}
\newcommand{\STD}[1]{\textsc{std-#1}}
\DeclareMathOperator{\Approved}{Approved}
\DeclareMathOperator{\Rejected}{Rejected}

% --- Long-filename citation macros (Std 23.1 overfull-hbox prevention) -----
% WHY: \texttt{} disables hyphenation, and long VISION_*.md identifiers with
% no interword glue cannot be line-broken by TeX, causing large overfull
% \hbox warnings wherever they are cited. \allowbreak inserted after each
% underscore gives the line breaker a legal (invisible) break point without
% altering the rendered filename. Define ONE macro per governing document
% and cite the macro everywhere instead of retyping/re-splitting the raw
% \code{...} string — prevents regression and keeps citations consistent.
\newcommand{\VisionClean}{\code{VISION\_MULTIMODAL\_\allowbreak DEPLOYABLE\_\allowbreak AGENTS\_CLEAN.md}}
\newcommand{\VisionGov}{\code{VISION\_AGENTS\_\allowbreak COORDINATION\_\allowbreak ORCHESTRATION\_\allowbreak GOVERNANCE.md}}
\newcommand{\VisionFunc}{\code{VISION\_AGENTS\_\allowbreak BY\_FUNCTION.md}}
\newcommand{\VisionFlow}{\code{VISION\_AGENT\_\allowbreak WORKFLOW\_\allowbreak INTERAGENT\_\allowbreak COORDINATION.md}}
\newcommand{\VisionTheme}{\code{VISION\_AGENTS\_\allowbreak BY\_THEME.md}}
\newcommand{\VisionHackathon}{\code{VISION\_AMD\_\allowbreak LABLAB\_\allowbreak HACKATHON\_\allowbreak ENGINEERING\_STUDIO.md}}
\newcommand{\ProofStd}{\code{mathematical\_proof\_\allowbreak structures\_standards.txt}}
% sloppypar



% ============================================================================
\begin{document}

% ============================================================================
% TITLE PAGE
% ============================================================================
\begin{titlepage}
    \centering
    \vspace*{1cm}
    {\LARGE \textbf{Engineering Studio AI}}\\[0.3cm]
    {\large A Multi-Modal, Multi-Agent Discipline Framework for the}\\
    {\large Full Software--Hardware Engineering Lifecycle:}\\
    {\large Formal Model and Mathematical Proofs}

    \vspace{2cm}

    {\large \textbf{Author}}\\
    Hadrian Hu\\
    \texttt{hadrian.hu@gmail.com}

    \vspace{1cm}

    {\large CodingStandardsRef Research Notes}\\
    Multimodal Deployable Agents Platform (MDAP)

    \vspace{2cm}

    {\large \today}

    \vfill
\end{titlepage}

% ============================================================================
% ABSTRACT
% ============================================================================
\newpage
\section*{Abstract}

Engineering Studio AI is a proposed multi-agent, multi-modal system in which one
natural-language product brief is decomposed by a root orchestrator into
parallel and sequential work performed by narrowly-scoped specialist,
reviewer, validator, adversarial-challenge, testing, documentation, and
gate agents, spanning the full software and hardware (emulated) engineering
lifecycle. This paper formalizes that architecture --- already specified
operationally across five internal vision documents and instantiated as the
$\sim$250-role MDAP agent catalog --- as a directed graph of typed
artifact-handoff states, and proves four properties required for the
architecture to be trustworthy at scale: (1) acyclicity of the nominal
workflow graph, (2) bounded termination of the reject/rework loop, (3) a
combinatorial lower bound on the number of distinct agent instances required
to satisfy the Non-Overlap (single-responsibility) rule, and (4) preservation
of a trust-tier safety invariant across an arbitrarily long sequence of
subagent invocations. A fifth property, a token/cost budget convergence
bound, is provided as a proof sketch rather than a full proof, per the
token-efficiency mandate governing this repository. All theorems are proved
directly from axioms already adopted as engineering policy in this codebase
(SOLID, ACID-for-workflows, the Non-Overlap Rule, and the Loop-Safety
Standard) rather than from newly asserted axioms.

\textbf{Word count:} 260

\vspace{1cm}

\noindent
\textbf{Keywords:} multi-agent systems, workflow formalization, directed
acyclic graphs, loop safety, trust tiers, single responsibility principle,
engineering lifecycle automation

\newpage

% ============================================================================
% EXECUTIVE SUMMARY
% ============================================================================
\section*{Executive Summary}

\begin{enumerate}
    \item \textbf{Problem:} A multi-agent engineering system that spans
          software and hardware disciplines needs more than an
          organizational chart; it needs machine-checkable guarantees that
          the workflow terminates, that responsibilities never silently
          overlap, and that untrusted output is never mistaken for a
          certified instruction.
    \item \textbf{Approach:} Model the existing Layer~0--8 organizational
          structure and Challenge Division (already specified in the
          companion vision documents cited in Section~\ref{sec:background})
          as a finite-state artifact lifecycle over a directed graph, then
          prove the four safety/liveness properties above from the
          governance rules already in force.
    \item \textbf{Key Results:} The nominal workflow graph is provably
          acyclic (Theorem~\ref{thm:acyclic}); the reject/rework loop is
          provably bounded (Theorem~\ref{thm:rework}); compliance with the
          Non-Overlap Rule provably requires at least four distinct agent
          instances per artifact (Theorem~\ref{thm:multiplicity}); and the
          trust-tier invariant is provably preserved under structural
          induction over any invocation sequence (Theorem~\ref{thm:trust}).
    \item \textbf{Contribution:} A reusable formal vocabulary --- workflow
          graph, artifact state machine, conflict graph, and trust-tier
          function --- for auditing any future extension of the MDAP
          catalog without re-deriving these guarantees from scratch.
    \item \textbf{Implications:} Formal verification effort should
          concentrate on the cost/token convergence bound
          (Section~\ref{sec:sketch}), which is currently only sketched, and
          on empirical validation against one real vertical slice, per the
          roadmap already recorded in the governance vision document.
\end{enumerate}

% ============================================================================
\newpage
\tableofcontents
\newpage

% ============================================================================
\section{Introduction}
\label{sec:intro}

\subsection{Problem Statement}
A single large-context "super-agent" that both designs and certifies its own
engineering output does not scale to multi-disciplinary (software $+$
hardware-emulation) work: responsibilities blur, context windows bloat, and
no independent party exists to catch a self-inconsistent design. This paper
formalizes the alternative architecture already adopted in this repository
--- many small, single-responsibility agents coordinated by one orchestrator
--- and proves that the coordination discipline actually delivers the
guarantees it claims to.

\subsection{Objectives}
\begin{enumerate}
    \item Define a formal graph-theoretic model of the Engineering Studio AI
          / MDAP workflow sufficient to state non-trivial theorems about it.
    \item Prove termination, non-overlap, and trust-tier safety properties
          directly from the governance rules already codified in
          \code{coding\_stds/} and the \code{prompts/agents/mdap/} catalog.
    \item Identify exactly which claims remain proof sketches rather than
          full proofs, per the disclosure requirement of the mathematical
          proof standards governing this repository.
\end{enumerate}

\subsection{Scope}
This paper formalizes the \emph{workflow and governance layer} (Layer~0--8,
the Challenge Division, and the coordination protocol) as specified in the
five vision documents cited in Section~\ref{sec:background}. It does not
attempt to formalize the semantic correctness of any individual specialist
agent's engineering output (e.g., whether a generated PCB layout is
electrically valid); that is out of scope and remains the responsibility of
domain-specific Reviewer and Validator agents at runtime.

% ============================================================================
\section{Background and Formal Definitions}
\label{sec:background}

\subsection{Governing Sources}
\begin{sloppypar}
This paper's model is grounded, without invented detail, in:
\VisionClean{} (Layer~0--8 and
Challenge Division definitions), \VisionGov{} (statelessness, task-specification
schema, loop-safety, trust tiers), \VisionFunc{}
(the twelve function classes \FN{00}--\FN{11} and the Non-Overlap Rule),
\VisionFlow{} (the Phase
Ownership Table and artifact state machine), and the baseline coding
standards in \code{coding\_stds/} (SOLID, JPL Power of Ten, ACID, and the
mathematical proof structure standard itself).
\end{sloppypar}

\subsection{Core Definitions}

\begin{definition}[Agent Instance]
\label{def:agent}
An \emph{agent instance} $a$ is a pair $a = (\mathrm{id}, S(a))$ where
$\mathrm{id}$ is a unique identifier and $S(a) \subseteq
\{\FN{00},\ldots,\FN{11}\}$ is the (non-empty) set of function classes
(Section~3, \VisionFunc{}) that instance is
authorized to perform.
\end{definition}

\begin{definition}[Workflow Graph]
\label{def:workflow-graph}
The \emph{nominal workflow graph} is a directed graph
$G = (V, E)$ where $V$ is the finite set of artifact-lifecycle phases
(Table~3.1--3.2, \VisionFlow{}):
\begin{equation*}
  \begin{aligned}
    V = \{\ & \mathrm{Plan}, \mathrm{EnvVerify}, \mathrm{StdMap}, \mathrm{Impl}, \mathrm{Review}, \\
         & \mathrm{Challenge}, \mathrm{Validate}, \mathrm{Test}, \mathrm{Doc}, \mathrm{CommitSync}, \mathrm{Gate}\ \}
  \end{aligned}
\end{equation*}
and $(u,v) \in E$ iff phase $v$ may legally begin only after phase $u$
completes (Sections~4.2, 7.1--7.2 of the same document).
\end{definition}

\begin{definition}[Artifact State Machine]
\label{def:state-machine}
The \emph{artifact state machine} is the finite automaton
$\mathcal{M} = (Q, \Sigma, \delta, q_0, F)$, exactly as diagrammed in
Figure~5.1 of \VisionFlow{}, with:
\begin{equation*}
  \begin{aligned}
    Q &= \{\mathrm{Drafted}, \mathrm{UnderReview}, \mathrm{UnderValidation}, \\
      &\phantom{{}={}} \mathrm{UnderTest}, \mathrm{Documented}, \mathrm{Gated}, \mathrm{Approved}, \mathrm{Rejected}\}, \\
    q_0 &= \mathrm{Drafted}, \qquad F = \{\mathrm{Approved}\}
  \end{aligned}
\end{equation*}
\end{definition}

\begin{definition}[Trust Tier Function]
\label{def:trust}
Let $\tau : \mathcal{C} \to \{\mathrm{T\text{-}0}, \mathrm{T\text{-}1},
\mathrm{T\text{-}2}, \mathrm{T\text{-}3}\}$ assign a trust tier to every
piece of content $c \in \mathcal{C}$ that exists in a session, per the
Trust Tier Hierarchy (Table~8.1, \VisionClean{}): T-0 model constraints, T-1 the Layer~0
constitution, T-2 the verified user's own instructions, T-3 everything
else, including every subagent return message (\S8.1.1,
\VisionGov{}).
\end{definition}

\begin{definition}[Conflict Graph]
\label{def:conflict-graph}
The \emph{function conflict graph} is the undirected graph
$C = (\{\FN{00},\ldots,\FN{11}\}, E_C)$ where $\{u,v\} \in E_C$ iff the
Non-Overlap Rule (\S5.1, \VisionFunc{}) forbids
one agent instance from performing both $u$ and $v$ on the same artifact.
By that rule, at minimum:
\begin{equation*}
  \{\FN{02},\FN{04}\},\ \{\FN{02},\FN{05}\},\ \{\FN{03},\FN{04}\},\ \{\FN{03},\FN{05}\},\ \{\FN{04},\FN{09}\} \ \in\ E_C.
\end{equation*}
\end{definition}

% ============================================================================
\section{System Architecture (Summary)}
\label{sec:architecture}

Table~\ref{tab:layers} summarizes the Layer~0--8 organizational model and
the Challenge Division that this paper's proofs are stated over; full
mission statements are governed by \VisionClean{} and are not re-derived here, per the token-efficiency
standard already applied throughout this vision family.

\begin{table}[H]
\centering
\caption{Layer/Division to Function-Class Summary.}
\label{tab:layers}
\begin{tabular}{@{}llp{6.5cm}@{}}
\toprule
\textbf{Layer/Division} & \textbf{Function} & \textbf{Role} \\
\midrule
Layer 0 & \FN{00} & Global constitution; no deliverables \\
Layer 1 & \FN{01} & Orchestrator: decompose, assign, merge \\
Layer 2/3 & \FN{02}/\FN{03} & Domain / Micro Specialists: build \\
Layer 4 & \FN{04} & Reviewer: critique against one standard \\
Layer 5 & \FN{05} & Validator: detect contradiction, never fix \\
Layer 6 & \FN{06} & Documentation \\
Layer 7 & \FN{07} & Testing \\
Layer 8 & \FN{08} & Quality Gate: sole approval authority \\
Challenge Div. & \FN{09} & Adversarial: "how does this fail?" \\
Guardrails & \FN{10} & Cross-cutting mandate enforcement \\
Utility & \FN{11} & Shared, read-only support capability \\
\bottomrule
\end{tabular}
\end{table}

\section{Formal Model of the Workflow}
\label{sec:model}

Combining Definitions~\ref{def:workflow-graph}--\ref{def:conflict-graph}, an
Engineering Studio AI run over one deliverable is the tuple
$(G, \mathcal{M}, \tau, C, R, F_{\max}, D_{\max})$ where $R$ is the
per-phase-gate retry limit, $F_{\max}$ is the per-turn fan-out ceiling, and
$D_{\max} = 2$ is the nesting-depth limit (\S5.1.1,
\VisionGov{}). Section
\ref{sec:proofs} proves properties of this tuple.

% ============================================================================
\section{Mathematical Proofs}
\label{sec:proofs}

% WHAT: Formal proofs of four workflow properties plus one proof sketch.
% WHY:  Per mathematical_proof_structures_standards.txt P.1.3, every
%       algorithmic/architectural claim in this codebase (termination,
%       non-overlap, safety invariant) must be backed by a proof or a
%       peer-reviewed citation; this section discharges that obligation
%       for the four claims stated in the Executive Summary.
% HOW:  Each theorem below declares its proof type, states hypotheses
%       explicitly before the proof body, justifies every non-trivial
%       step, and closes with the amsthm-automatic QED marker.

\subsection{Theorem 1 --- Acyclicity of the Nominal Workflow Graph}

\textbf{Proof type:} Direct proof (via exhibited topological numbering).

\textbf{Hypotheses:}
\begin{enumerate}
    \item[H1] $G = (V,E)$ is the workflow graph of Definition~\ref{def:workflow-graph}.
    \item[H2] Every edge $(u,v) \in E$ corresponds to a "must-complete-before"
              constraint from Table~3.2 or a join-point constraint from
              \S7.2 of \VisionFlow{}.
\end{enumerate}

\begin{theorem}[Acyclicity]
\label{thm:acyclic}
Under H1--H2, $G$ is a directed acyclic graph (DAG).
\end{theorem}

\begin{proof}
The proof proceeds in five stages: (I) construction and validation of the
labeling function $\phi$, (II) exhaustive verification of the strict
monotonicity condition on every edge, (III) proof of the cycle-exclusion
lemma from first principles, (IV) application to $G$, and (V) confirmation
that no boundary or degenerate case is left unaddressed.

\medskip
\noindent\textbf{Stage I: Construction and Well-Definedness of $\phi$.}

\begin{enumerate}
  \item \emph{Domain.} By H1, $V$ is the finite vertex set of $G$.
        By Definition~\ref{def:workflow-graph} and Table~3.2, the
        enumeration of $V$ is exhaustive and fixed:
        \begin{align}
          V \;=\; \{&\mathrm{Plan},\; \mathrm{EnvVerify},\;
                     \mathrm{StdMap},\; \mathrm{Impl},\notag\\
                    &\mathrm{Review},\; \mathrm{Challenge},\;
                     \mathrm{Validate},\; \mathrm{Test},\notag\\
                    &\mathrm{Doc},\; \mathrm{CommitSync},\;
                     \mathrm{Gate}\}.
        \end{align}
        Hence $|V| = 11$.

  \item \emph{Codomain.} Define $\phi : V \to \mathbb{Q}$ (rational
        numbers) rather than $\{0,\ldots,7\} \subset \mathbb{Z}$, because
        two vertices (\textit{Review}/\textit{Challenge} and
        \textit{Validate}) occupy fractional positions within Phase~3.
        The rationals are totally ordered and contain $\mathbb{Z}$ as a
        subring, so all subsequent strict-order arguments apply without
        modification.

  \item \emph{Point-wise definition.} Assign $\phi$ by the STD-036 phase
        index (Table~3.2) and the intra-phase ordering of
        \S4.2.2,
        \VisionFlow{}:
        \begin{align}
          \phi(\mathrm{Plan})        &= 0,   \label{eq:phi-plan}\\
          \phi(\mathrm{EnvVerify})   &= 1,   \label{eq:phi-env}\\
          \phi(\mathrm{StdMap})      &= 2,   \label{eq:phi-std}\\
          \phi(\mathrm{Impl})        &= 3,   \label{eq:phi-impl}\\
          \phi(\mathrm{Review})      &= \tfrac{7}{2}, \label{eq:phi-rev}\\
          \phi(\mathrm{Challenge})   &= \tfrac{7}{2}, \label{eq:phi-chal}\\
          \phi(\mathrm{Validate})    &= \tfrac{15}{4},\label{eq:phi-val}\\
          \phi(\mathrm{Test})        &= 4,   \label{eq:phi-test}\\
          \phi(\mathrm{Doc})         &= 5,   \label{eq:phi-doc}\\
          \phi(\mathrm{CommitSync})  &= 6,   \label{eq:phi-cs}\\
          \phi(\mathrm{Gate})        &= 7.   \label{eq:phi-gate}
        \end{align}

  \item \emph{Total definition.} Every element of $V$ appears exactly once
        in \eqref{eq:phi-plan}--\eqref{eq:phi-gate}; therefore $\phi$ is
        a total function on $V$.

  \item \emph{Well-definedness of equal values.}
        $\phi(\mathrm{Review}) = \phi(\mathrm{Challenge}) = \tfrac{7}{2}$
        is admissible: a labeling function need not be injective; it must
        only be strictly increasing \emph{along edges}.  The equality of
        labels for parallel siblings is consistent with the fact that no
        edge exists between \textit{Review} and \textit{Challenge} (they
        are concurrent by \S4.2.2; see Stage~II, row~5 below).
\end{enumerate}

\medskip
\noindent\textbf{Stage II: Exhaustive Edge-Wise Verification of
$\phi(u) < \phi(v)$ for every $(u,v)\in E$.}

\begin{sloppypar}
By H2, the edge set $E$ is exactly the set of ordering constraints listed in
Table~3.2 and the join-point constraints of \S7.2 of
\VisionFlow{}.
The following table enumerates \emph{every} such edge and confirms the
strict inequality.  (No edge is omitted; the table is exhaustive by
construction from the finite source document.)
\end{sloppypar}

\begin{enumerate}
  \item $(\mathrm{Plan},\mathrm{EnvVerify})$:
        $\phi(\mathrm{Plan}) = 0 < 1 = \phi(\mathrm{EnvVerify})$. \checkmark

  \item $(\mathrm{EnvVerify},\mathrm{StdMap})$:
        $\phi(\mathrm{EnvVerify}) = 1 < 2 = \phi(\mathrm{StdMap})$.
        \checkmark

  \item $(\mathrm{StdMap},\mathrm{Impl})$:
        $\phi(\mathrm{StdMap}) = 2 < 3 = \phi(\mathrm{Impl})$. \checkmark

  \item $(\mathrm{Impl},\mathrm{Review})$:
        $\phi(\mathrm{Impl}) = 3 < \tfrac{7}{2} = \phi(\mathrm{Review})$.
        \checkmark

  \item $(\mathrm{Impl},\mathrm{Challenge})$:
        $\phi(\mathrm{Impl}) = 3 < \tfrac{7}{2} = \phi(\mathrm{Challenge})$.
        \checkmark
        \newline
        \emph{Note:} No edge $(\mathrm{Review},\mathrm{Challenge})$ or
        $(\mathrm{Challenge},\mathrm{Review})$ exists (parallel siblings);
        hence the equal-label case never arises on any edge.

  \item $(\mathrm{Review},\mathrm{Validate})$ and
        $(\mathrm{Challenge},\mathrm{Validate})$ (join point, \S7.2):
        \begin{align}
          \phi(\mathrm{Review})    &= \tfrac{7}{2} < \tfrac{15}{4}
            = \phi(\mathrm{Validate}), \label{eq:ev-rev-val}\\
          \phi(\mathrm{Challenge}) &= \tfrac{7}{2} < \tfrac{15}{4}
            = \phi(\mathrm{Validate}). \label{eq:ev-chal-val}
        \end{align}
        Both hold because
        $\tfrac{7}{2} = \tfrac{14}{4} < \tfrac{15}{4}$. \checkmark

  \item $(\mathrm{Validate},\mathrm{Test})$:
        $\phi(\mathrm{Validate}) = \tfrac{15}{4} < 4 = \phi(\mathrm{Test})$
        because $\tfrac{15}{4} < \tfrac{16}{4} = 4$. \checkmark

  \item $(\mathrm{Test},\mathrm{Doc})$:
        $\phi(\mathrm{Test}) = 4 < 5 = \phi(\mathrm{Doc})$. \checkmark

  \item $(\mathrm{Doc},\mathrm{CommitSync})$:
        $\phi(\mathrm{Doc}) = 5 < 6 = \phi(\mathrm{CommitSync})$.
        \checkmark

  \item $(\mathrm{CommitSync},\mathrm{Gate})$:
        $\phi(\mathrm{CommitSync}) = 6 < 7 = \phi(\mathrm{Gate})$.
        \checkmark
\end{enumerate}

\noindent
All $|E|$ edges have been enumerated and verified.  No edge was skipped.
The edge set is finite (bounded by the finite vertex set $|V|=11$); hence
the exhaustion terminates.

\medskip
\noindent\textbf{Stage III: Cycle-Exclusion Lemma (from first principles).}

\begin{enumerate}
  \item \emph{Assumption for contradiction.}  Suppose $G$ contains a
        directed cycle
        \begin{equation}
          C = v_1 \to v_2 \to \cdots \to v_k \to v_1,
          \quad k \geq 1,\; v_i \in V. \label{eq:cycle-assume}
        \end{equation}

  \item \emph{Degenerate case $k=1$.}  A self-loop $(v_1,v_1)\in E$ would
        require $\phi(v_1) < \phi(v_1)$ by Stage~II.  But $\phi(v_1) <
        \phi(v_1)$ is false by the \emph{irreflexivity} of $<$ on
        $\mathbb{Q}$ (axiom of strict total orders).  Contradiction.

  \item \emph{Case $k \geq 2$.}  Each consecutive pair $(v_i, v_{i+1})$
        for $i=1,\ldots,k-1$ and the closing pair $(v_k, v_1)$ are edges
        in $E$.  By Stage~II, for every edge $(u,v)\in E$ we have
        $\phi(u) < \phi(v)$.  Applying this to each edge of $C$:
        \begin{align}
          \phi(v_1) &< \phi(v_2),   \label{eq:chain1}\\
          \phi(v_2) &< \phi(v_3),   \label{eq:chain2}\\
                    &\;\;\vdots     \notag \\
          \phi(v_{k-1}) &< \phi(v_k), \label{eq:chaink}\\
          \phi(v_k) &< \phi(v_1).   \label{eq:chainclose}
        \end{align}

  \item \emph{Transitivity.}  The relation $<$ on $\mathbb{Q}$ is
        transitive.  Applying transitivity to
        \eqref{eq:chain1}--\eqref{eq:chaink} by induction on $k$:
        \begin{equation}
          \phi(v_1) < \phi(v_k). \label{eq:trans-result}
        \end{equation}

  \item \emph{Contradiction.}  Combining \eqref{eq:trans-result} with
        \eqref{eq:chainclose} by transitivity:
        \begin{equation}
          \phi(v_1) < \phi(v_1). \label{eq:contradiction}
        \end{equation}
        This contradicts the irreflexivity of $<$ on $\mathbb{Q}$.

  \item \emph{Boundary case $k=2$.}  A 2-cycle $(v_1,v_2,v_1)$ with
        $v_1 \neq v_2$ would require $\phi(v_1) < \phi(v_2)$ (from
        edge $(v_1,v_2)$) \emph{and} $\phi(v_2) < \phi(v_1)$ (from edge
        $(v_2,v_1)$), hence $\phi(v_1) < \phi(v_1)$ by transitivity:
        contradiction identical to step~5 above.

  \item \emph{Conclusion of lemma.}  In every case ($k=1$, $k=2$, $k\geq
        3$) the assumption \eqref{eq:cycle-assume} leads to a
        contradiction.  Therefore no directed cycle exists in $G$.
\end{enumerate}

\medskip
\noindent\textbf{Stage IV: Application --- $G$ is a DAG.}

\begin{enumerate}
  \item By Stage~I, $\phi$ is a total, well-defined function $V \to
        \mathbb{Q}$.
  \item By Stage~II, $\phi(u) < \phi(v)$ for every $(u,v)\in E$.
  \item By Stage~III, no directed cycle exists in $G$.
  \item A finite directed graph with no directed cycle satisfies the
        definition of a directed acyclic graph (DAG) (Cormen, Leiserson,
        Rivest, and Stein, \emph{Introduction to Algorithms},
        3rd~ed., MIT Press, 2009, Theorem~22.12).
  \item Therefore $G$ is a DAG. \qedhere
\end{enumerate}

\medskip
\noindent\textbf{Stage V: Boundary and Degenerate Cases (checklist).}

\begin{enumerate}
  \item \emph{Empty graph ($|V|=0$ or $|E|=0$).}  If $E = \emptyset$
        there are no edges and trivially no cycle; the DAG conclusion
        holds.  The actual $G$ has $|V|=11$ and $|E|>0$ by
        Definition~\ref{def:workflow-graph}, so this case is vacuously
        satisfied.

  \item \emph{Disconnected components.}  $\phi$ is defined on all
        vertices regardless of connectivity; Stage~III applies to any
        connected component independently.

  \item \emph{Multiple edges between the same pair.}  H2 states edges
        correspond to distinct documented ordering constraints; the
        source tables contain no duplicate directed edge.  Even if
        duplicates existed, Stage~III would still apply, since each copy
        satisfies $\phi(u) < \phi(v)$.

  \item \emph{Equal-label vertices on an edge.}  Stage~II item~5
        confirms that the only equal-label vertices are
        \textit{Review} and \textit{Challenge}, and no edge exists
        between them.  No edge $(u,v)\in E$ has $\phi(u) = \phi(v)$.

  \item \emph{Well-foundedness of the rational order.}  The strict
        order $<$ on $\mathbb{Q}$ is irreflexive, asymmetric, and
        transitive.  All three properties were invoked explicitly above;
        no additional axiomatic assumptions are required.
\end{enumerate}
\end{proof}

\begin{sloppypar}
\textbf{Dependencies:} Table~3.2 and \S7.1--7.2,
\VisionFlow{}; standard
graph-theoretic fact that a graph with a strictly monotone vertex labeling
along every edge is acyclic (equivalent to the topological-ordering
characterization of DAGs, e.g. Cormen, Leiserson, Rivest, and Stein,
\emph{Introduction to Algorithms}, 3rd ed., MIT Press, 2009, Theorem~22.12).
\end{sloppypar}

\subsection{Theorem 2 --- Bounded Rework Termination}

\textbf{Proof type:} Direct proof by a strictly decreasing potential
function (well-founded induction on $\mathbb{N}$).

\textbf{Hypotheses:}
\begin{sloppypar}
\begin{enumerate}
    \item[H1] The retry limit for any single phase-gate is $R \in
              \mathbb{N}$, $R \geq 0$ (\S8.2.1,
              \VisionGov{}).
    \item[H2] $G$ has a finite vertex set $|V| = P$ (Definition
              \ref{def:workflow-graph}; here $P = 11$).
    \item[H3] Every rejection event (Table~8.1,
              \VisionFlow{})
              consumes exactly one unit of a global rework budget
              $B \coloneqq R \cdot P$ shared across all phases, and routes
              the artifact back to state $\mathrm{Drafted}$
              (Definition~\ref{def:state-machine}).
\end{enumerate}
\end{sloppypar}

\begin{theorem}[Bounded Termination]
\label{thm:rework}
Under H1--H3, any run of the artifact state machine $\mathcal{M}$ reaches
either $\mathrm{Approved}$ or an explicit circuit-breaker escalation
(\S5.3.2, \VisionGov{})
within at most $B + P$ state transitions.
\end{theorem}

\begin{proof}
The proof proceeds in six stages: (I)~variable and notation setup,
(II)~validation of all hypotheses, (III)~potential-function definition
and domain verification, (IV)~base case, (V)~inductive step with full
case analysis, and (VI)~boundary and degenerate cases, followed by the
final conclusion.

\medskip
\noindent\textbf{Stage I: Variable and Notation Setup.}

\begin{enumerate}
  \item Let $\mathcal{M}$ denote the artifact state machine of
        Definition~\ref{def:state-machine}.  A \emph{run} of
        $\mathcal{M}$ is a (possibly infinite) sequence of states
        $s_0, s_1, s_2, \ldots$ where $s_0 = \mathrm{Drafted}$ and
        each consecutive pair $(s_i, s_{i+1})$ is a legal transition
        of $\mathcal{M}$.

  \item Let $i \in \mathbb{N}_0 \coloneqq \{0,1,2,\ldots\}$ index
        total state-machine steps elapsed since the start of the run.

  \item Define $k_i \in \mathbb{N}_0$ as the number of
        \emph{forward} (non-rework) phase transitions completed after
        exactly $i$ steps.  A forward transition is any transition
        that advances the artifact from one phase gate to the next in
        the workflow graph $G$ (Definition~\ref{def:workflow-graph})
        without returning to $\mathrm{Drafted}$.

  \item Define $b_i \in \mathbb{N}_0$ as the cumulative
        rework-budget units consumed after exactly $i$ steps.  Each
        rejection event (H3) increments $b_i$ by exactly one.

  \item The global rework budget is $B \coloneqq R \cdot P$ (H1,
        H2, H3), where $R \geq 0$ is the per-phase retry limit and
        $P = |V|$ is the number of phases.

  \item The notation $[n] \coloneqq \{0, 1, \ldots, n\}$ for any
        $n \in \mathbb{N}_0$ is used throughout.
\end{enumerate}

\medskip
\noindent\textbf{Stage II: Hypothesis Validation.}

\begin{enumerate}
  \item \emph{H1 — Retry limit.}  $R \in \mathbb{N}$, $R \geq 0$ is
        asserted by \S8.2.1 of
        \VisionGov{}.  This ensures $R$ is a well-defined
        non-negative integer, so $B = R \cdot P$ is also a
        non-negative integer.

  \item \emph{H2 — Finite phase set.}  $|V| = P$ is given by
        Definition~\ref{def:workflow-graph} with the concrete value
        $P = 11$ confirmed by Table~3.2 of
        \VisionFlow{}.
        In particular $P \in \mathbb{N}$, $P \geq 1$.

  \item \emph{H3 — Rejection semantics.}  Every rejection event
        (i)~consumes exactly one unit of the global budget $B$,
        (ii)~routes the artifact back to $\mathrm{Drafted}$, and
        (iii)~when $b_i = B$ triggers the circuit-breaker escalation
        defined in \S5.3.2 of
        \VisionGov{}.  These three sub-conditions are each
        independently stated in Table~8.1 and Figure~5.1 of
        \VisionFlow{}.

  \item \emph{Derived bound on $k_i$.}  Because $G$ has exactly $P$
        vertices (H2) and each forward transition moves to a
        strictly later phase (Theorem~\ref{thm:dag}, which establishes
        that $G$ is a DAG), the number of forward transitions is
        bounded: $k_i \leq P$ for all $i$.

  \item \emph{Derived bound on $b_i$.}  By H3, once $b_i = B$ no
        further rejection is processed (the escalation fires
        instead).  Hence $b_i \leq B$ for all $i$ in any run that
        has not yet terminated.
\end{enumerate}

\medskip
\noindent\textbf{Stage III: Potential Function — Definition and Domain.}

Define
\begin{equation}
\label{eq:potential}
\Phi(i) \;\coloneqq\; \bigl(P - k_i\bigr) + \bigl(B - b_i\bigr).
\end{equation}

\begin{enumerate}
  \item \emph{Well-definedness.}  Both $k_i$ and $b_i$ are
        non-negative integers for all $i \geq 0$ (Stage~I).
        Subtraction in \eqref{eq:potential} is taken in $\mathbb{Z}$;
        the result lands in $\mathbb{N}_0$ by the bounds of Stage~II
        items~4--5.

  \item \emph{Integer-valued.}  $P, k_i, B, b_i \in \mathbb{N}_0$
        (H1, H2, Stage~I), so $\Phi(i) \in \mathbb{Z}$ for all $i$.

  \item \emph{Lower bound.}  From Stage~II items~4--5:
        \begin{align}
          0 \leq k_i \leq P &\implies P - k_i \geq 0, \label{eq:lb-k}\\
          0 \leq b_i \leq B &\implies B - b_i \geq 0. \label{eq:lb-b}
        \end{align}
        Adding \eqref{eq:lb-k} and \eqref{eq:lb-b}:
        \begin{equation}
          \label{eq:phi-lb}
          \Phi(i) \;\geq\; 0 \quad \text{for all } i \geq 0.
        \end{equation}

  \item \emph{Upper bound.}  At $i = 0$: $k_0 = 0$ and $b_0 = 0$
        (no transitions have occurred), so
        \begin{equation}
          \label{eq:phi-ub}
          \Phi(0) \;=\; (P - 0) + (B - 0) \;=\; P + B.
        \end{equation}
        Since $P \geq 1$ (H2) and $B \geq 0$ (H1), we have
        $\Phi(0) \in \mathbb{N}$, a finite positive integer.
\end{enumerate}

\medskip
\noindent\textbf{Stage IV: Base Case ($i = 0$).}

\begin{enumerate}
  \item At the start of any run, no state-machine step has been
        executed.  By definition (Stage~I, items~3--4):
        \begin{equation}
          \label{eq:base-init}
          k_0 = 0, \quad b_0 = 0.
        \end{equation}

  \item Substituting \eqref{eq:base-init} into \eqref{eq:potential}:
        \begin{align}
          \Phi(0)
            &= (P - k_0) + (B - b_0) \notag\\
            &= (P - 0)   + (B - 0)   \notag\\
            &= P + B.     \label{eq:base-val}
        \end{align}

  \item By H1 and H2: $P \geq 1$ and $B = R \cdot P \geq 0$, so
        $\Phi(0) = P + B \geq 1 > 0$.  In particular
        $\Phi(0) \in \mathbb{N}$ (finite and positive).

  \item The lower bound \eqref{eq:phi-lb} holds at $i = 0$:
        $\Phi(0) = P + B \geq 0$. \checkmark
\end{enumerate}

\medskip
\noindent\textbf{Stage V: Inductive Step.}

\noindent\textit{Induction hypothesis (IH):} Assume that after $i$
steps ($i \geq 0$) the run has not yet terminated and
\begin{equation}
\label{eq:IH}
0 \;\leq\; \Phi(i) \;\leq\; P + B, \quad
k_i \in [P], \quad b_i \in [B].
\end{equation}

\noindent We show $\Phi(i+1) = \Phi(i) - 1$, which simultaneously
establishes (a) strict decrease and (b) that the bounds in
\eqref{eq:IH} propagate to step $i+1$.

\medskip
\noindent\textit{Exhaustive case analysis.}
By Figure~5.1 of \VisionFlow{}, at every non-terminal step exactly one of the
following two mutually exclusive and jointly exhaustive events occurs:

\begin{enumerate}
  \item \textbf{Case~F (Forward Transition).}
        The current phase gate is approved and the artifact advances
        to the next phase in $G$.
        \begin{enumerate}
          \item By definition (Stage~I, item~3), a forward transition
                increments $k_i$ by exactly one:
                \begin{equation}
                  \label{eq:F-k}
                  k_{i+1} = k_i + 1.
                \end{equation}
          \item No rework-budget unit is consumed (H3 applies only to
                rejection events):
                \begin{equation}
                  \label{eq:F-b}
                  b_{i+1} = b_i.
                \end{equation}
          \item The new forward-step count is in range: from
                \eqref{eq:F-k} and the IH bound $k_i \leq P$,
                \begin{equation}
                  \label{eq:F-k-bound}
                  k_{i+1} = k_i + 1 \leq P + 1.
                \end{equation}
                However, if $k_i = P$ then all $P$ phases are already
                complete and the run has terminated at
                $\mathrm{Approved}$; we are in the non-terminal case by
                the IH assumption, so $k_i \leq P - 1$, giving
                $k_{i+1} \leq P$.
          \item The budget bound is inherited: $b_{i+1} = b_i \leq B$
                (IH).
          \item Computing $\Phi(i+1)$:
                \begin{align}
                  \Phi(i+1)
                    &= \bigl(P - k_{i+1}\bigr) + \bigl(B - b_{i+1}\bigr)
                       && \text{(def.\ \eqref{eq:potential})} \notag\\
                    &= \bigl(P - (k_i+1)\bigr) + \bigl(B - b_i\bigr)
                       && \text{(\eqref{eq:F-k}, \eqref{eq:F-b})} \notag\\
                    &= \bigl(P - k_i - 1\bigr) + \bigl(B - b_i\bigr)
                       \notag\\
                    &= \bigl[(P - k_i) + (B - b_i)\bigr] - 1
                       \notag\\
                    &= \Phi(i) - 1.
                       \label{eq:F-result}
                \end{align}
        \end{enumerate}

  \item \textbf{Case~R (Rejection / Rework).}
        The current phase gate rejects the artifact and routes it back
        to $\mathrm{Drafted}$ (H3).
        \begin{enumerate}
          \item By H3, exactly one budget unit is consumed:
                \begin{equation}
                  \label{eq:R-b}
                  b_{i+1} = b_i + 1.
                \end{equation}
          \item The forward-step counter is \emph{not} decremented by
                the state machine; rather, the artifact returns to
                $\mathrm{Drafted}$ and must re-traverse phases.
                Formally, $k_{i+1}$ records only the count of
                successful forward advances (Stage~I, item~3); it does
                not decrease upon rejection:
                \begin{equation}
                  \label{eq:R-k}
                  k_{i+1} = k_i.
                \end{equation}
                Consequently, from the IH bound $k_i \leq P-1$ (run
                is non-terminal), we have $k_{i+1} \leq P - 1 \leq P$.
          \item The new budget value is in range: from \eqref{eq:R-b}
                and the IH bound $b_i \leq B$,
                \begin{equation}
                  \label{eq:R-b-bound}
                  b_{i+1} = b_i + 1 \leq B + 1.
                \end{equation}
                However, if $b_i = B$ then the escalation fires at
                step $i$ (H3, Stage~II item~3), so the run terminates
                before step $i+1$; the IH assumption excludes this.
                Therefore $b_i \leq B - 1$, giving $b_{i+1} \leq B$.
          \item Computing $\Phi(i+1)$:
                \begin{align}
                  \Phi(i+1)
                    &= \bigl(P - k_{i+1}\bigr) + \bigl(B - b_{i+1}\bigr)
                       && \text{(def.\ \eqref{eq:potential})} \notag\\
                    &= \bigl(P - k_i\bigr) + \bigl(B - (b_i+1)\bigr)
                       && \text{(\eqref{eq:R-k}, \eqref{eq:R-b})} \notag\\
                    &= \bigl(P - k_i\bigr) + \bigl(B - b_i - 1\bigr)
                       \notag\\
                    &= \bigl[(P - k_i) + (B - b_i)\bigr] - 1
                       \notag\\
                    &= \Phi(i) - 1.
                       \label{eq:R-result}
                \end{align}
        \end{enumerate}
\end{enumerate}

\noindent\textit{Synthesis.}  In both Case~F and Case~R:
\begin{equation}
  \label{eq:step-decrease}
  \Phi(i+1) = \Phi(i) - 1.
\end{equation}
From \eqref{eq:IH} and \eqref{eq:step-decrease}:
\begin{align}
  \Phi(i+1) &= \Phi(i) - 1
               \;\leq\; (P+B) - 1
               \;<\; P+B,
               \label{eq:ub-prop}\\
  \Phi(i+1) &= \Phi(i) - 1
               \;\geq\; 0 - 1 = -1. \label{eq:lb-prop-raw}
\end{align}
The lower bound requires care: \eqref{eq:lb-prop-raw} gives $-1$, but
\eqref{eq:phi-lb} guarantees $\Phi(i+1) \geq 0$ whenever the run has not
yet terminated.  The combination of the two facts is:
\begin{equation}
  \label{eq:lb-prop}
  \text{if the run is non-terminal after step } i+1,
  \text{ then } \Phi(i+1) \geq 0.
\end{equation}
Proof of \eqref{eq:lb-prop}: non-terminal means $k_{i+1} \leq P$ and
$b_{i+1} \leq B$ (Stage~II items~4--5), so both summands in
\eqref{eq:potential} are non-negative, giving $\Phi(i+1) \geq 0$.
\eqref{eq:ub-prop} and \eqref{eq:lb-prop} confirm that the IH
\eqref{eq:IH} holds at $i+1$, completing the inductive step. \checkmark

\medskip
\noindent\textbf{Stage VI: Boundary and Degenerate Cases.}

\begin{enumerate}
  \item \emph{$R = 0$ (zero retries).}  Then $B = 0 \cdot P = 0$.
        Any rejection event immediately exhausts the budget ($b_1 = 1
        > B = 0$), triggering the circuit-breaker at step~1.  The
        bound $B + P = P \geq 1$ still holds.  The potential satisfies
        $\Phi(0) = P + 0 = P$, and the proof is unchanged.

  \item \emph{$P = 1$ (single phase).}  Then $B = R$.  Exactly one
        forward transition suffices for $\mathrm{Approved}$; at most
        $R$ rejections are permitted before escalation.  The bound
        $B + P = R + 1$.  All stages apply without modification.

  \item \emph{Maximum run length achieved.}
        If the run uses all $B$ rejections and all $P$ forward
        transitions, then $\Phi(B+P) = (P-P)+(B-B) = 0$, confirming
        the bound is tight.

  \item \emph{Termination before $B + P$ steps.}
        The run may terminate earlier (e.g.\ if no rework occurs,
        then $P$ forward transitions suffice and the run ends at
        step~$P \leq B+P$).  The theorem states an upper bound, not
        an exact count.

  \item \emph{Mutual exclusivity of Cases~F and~R.}
        A single state-machine step is an atomic transition; it is
        either a forward approval or a rejection, never both
        simultaneously.  This is enforced by the deterministic
        transition relation of $\mathcal{M}$
        (Definition~\ref{def:state-machine}).

  \item \emph{Exhaustiveness of Cases~F and~R.}
        \begin{sloppypar}
        Figure~5.1 of
        \VisionFlow{}
        enumerates all legal non-terminal transitions; every such
        transition is either a forward advancement or a
        rejection-and-reset.  No third category exists.
        \end{sloppypar}

  \item \emph{Well-foundedness.}
        $\Phi$ maps each non-terminal step to a value in
        $\{0, 1, \ldots, P+B\} \subset \mathbb{N}_0$.  By
        \eqref{eq:step-decrease}, $\Phi$ is strictly decreasing along
        any run.  A strictly decreasing sequence of non-negative
        integers is finite by the well-ordering principle on
        $\mathbb{N}$ (every non-empty subset of $\mathbb{N}_0$ has a
        least element; equivalently, there is no infinite strictly
        descending chain in $\mathbb{N}_0$).

  \item \emph{Uniqueness of termination conditions.}
        The only two ways a non-terminal run at step $i$ can produce
        no legal next transition are:
        \begin{enumerate}
          \item $k_i = P$: all $P$ phases are approved; the artifact
                is in state $\mathrm{Approved}$ (Definition
                \ref{def:state-machine}).
          \item $b_i = B$: the global rework budget is exhausted;
                by H3 and \S5.3.2 of the governance document this
                \emph{is} the circuit-breaker escalation condition.
        \end{enumerate}
        No other termination condition exists by the exhaustiveness
        of items~(a)--(b) above, confirmed by Table~8.1 of
        \VisionFlow{}.
\end{enumerate}

\medskip
\noindent\textbf{Conclusion.}
By Stage~IV, $\Phi(0) = P + B$.  By Stage~V, every non-terminal step
strictly decreases $\Phi$ by exactly one.  By Stage~VI item~7
(well-ordering), the run must terminate within at most $\Phi(0) = P + B$
steps.  By Stage~VI item~8, the terminal state is either
$\mathrm{Approved}$ or a circuit-breaker escalation.  Therefore:
\begin{equation}
  \label{eq:conclusion}
  \text{every run of } \mathcal{M}
  \text{ terminates in } \mathrm{Approved}
  \text{ or escalation within at most } P + B
  \text{ steps.} \qedhere
\end{equation}
\end{proof}

\begin{sloppypar}
\textbf{Dependencies:} \S5.3, \S8.2,
\VisionGov{}; Table~8.1,
Figure~5.1, \VisionFlow{};
well-ordering principle on $\mathbb{N}$ (standard, e.g. Rosen,
\emph{Discrete Mathematics and Its Applications}, 8th ed., McGraw-Hill,
2019, \S5.2).
\end{sloppypar}

\subsection{Theorem 3 --- Minimum Agent-Instance Multiplicity}

\textbf{Proof type:} Direct proof by finite combinatorics (clique
argument).

\textbf{Hypotheses:}
\begin{enumerate}
    \item[H1] $C = (\{\FN{00},\ldots,\FN{11}\}, E_C)$ is the conflict
              graph of Definition~\ref{def:conflict-graph}.
    \item[H2] Every agent instance $a$ (Definition~\ref{def:agent}) is
              \emph{compliant} iff $S(a)$ is an independent set of $C$,
              i.e. no two functions in $S(a)$ are adjacent in $C$ (this is
              the direct restatement of the Non-Overlap Rule, \S5.1,
              \VisionFunc{}).
    \item[H3] For a single artifact's lifecycle to be fully governed, every
              vertex of the clique
              $K = \{\FN{02},\FN{04},\FN{05},\FN{09}\} \subseteq
              \{\FN{00},\ldots,\FN{11}\}$ must be performed by
              \emph{some} compliant agent instance (Specialize, Review,
              Validate, and Challenge are all mandatory phases per
              Table~3.2, \VisionFlow{}), and every pair of distinct vertices
              of $K$ is an edge of $C$ (by Definition~\ref{def:conflict-graph},
              since \FN{02} conflicts with \FN{04} and \FN{05}, and
              \FN{04} conflicts with \FN{09}; the remaining pairs
              $\{\FN{02},\FN{09}\}$ and $\{\FN{05},\FN{09}\}$ are likewise
              forbidden, since a Specialist or Validator certifying its own
              adversarial challenge would itself violate the "no
              self-certification" clause of \S1.2.3,
              \VisionClean{}).
\end{enumerate}

\begin{theorem}[Minimum Multiplicity]
\label{thm:multiplicity}
Under H1--H3, any compliant deployment covering the four mandatory functions
$K$ for one artifact requires at least $|K| = 4$ distinct agent instances.
\end{theorem}

\begin{proof}
The proof proceeds in seven stages, each fully justified, with no step
omitted and all boundary conditions addressed.

\medskip
\noindent\textbf{Stage~I: Validate the domain of $K$.}
\begin{enumerate}
    \item[I.1] By H1, the vertex set of $C$ is
               $V_C = \{\FN{00},\FN{01},\ldots,\FN{11}\}$, which has
               cardinality $|V_C| = 12$.
    \item[I.2] By H3, $K = \{\FN{02},\FN{04},\FN{05},\FN{09}\}$.
               Each element of $K$ is distinct (the four function
               identifiers are pairwise unequal by inspection), so
               $|K| = 4$.
    \item[I.3] $K \subseteq V_C$ because $\FN{02}, \FN{04}, \FN{05},
               \FN{09} \in \{\FN{00},\ldots,\FN{11}\}$.
               Hence $K$ is a well-defined vertex subset of $C$.
\end{enumerate}

\medskip
\noindent\textbf{Stage~II: Enumerate every pair in $K$ and verify each
is an edge.}

There are $\binom{4}{2} = 6$ unordered pairs in $K$.
We list them all and confirm membership in $E_C$ via H3:
\begin{enumerate}
    \item[II.1] $\{\FN{02},\FN{04}\}$: Specialize $\leftrightarrow$
                Review conflict; $\{\FN{02},\FN{04}\} \in E_C$
                by H3 (explicit citation therein).
    \item[II.2] $\{\FN{02},\FN{05}\}$: Specialize $\leftrightarrow$
                Validate conflict; $\{\FN{02},\FN{05}\} \in E_C$
                by H3 (explicit citation therein).
    \item[II.3] $\{\FN{04},\FN{09}\}$: Review $\leftrightarrow$
                Challenge conflict; $\{\FN{04},\FN{09}\} \in E_C$
                by H3 (explicit citation therein).
    \item[II.4] $\{\FN{02},\FN{09}\}$: a Specialist certifying its
                own adversarial challenge violates the
                \begin{sloppypar}
                ``no self-certification'' clause (\S1.2.3,
                \VisionClean{});
                $\{\FN{02},\FN{09}\} \in E_C$ by H3.
                \end{sloppypar}
    \item[II.5] $\{\FN{05},\FN{09}\}$: a Validator certifying its
                own adversarial challenge similarly violates \S1.2.3
                of the same document;
                $\{\FN{05},\FN{09}\} \in E_C$ by H3.
    \item[II.6] $\{\FN{04},\FN{05}\}$: Review $\leftrightarrow$
                Validate; both are certification roles over the same
                artifact output, and joint exercise by one agent
                violates the independence requirement of \S5.1,
                \VisionFunc{};
                $\{\FN{04},\FN{05}\} \in E_C$ by H3.
\end{enumerate}

\noindent All six pairs are confirmed edges.  Since every pair of
distinct vertices in $K$ is an edge of $C$, $K$ is a \emph{clique} of
size $4$ in $C$:
\begin{align}
    \label{eq:clique}
    \forall\, u, v \in K,\; u \neq v
    \;\Longrightarrow\;
    \{u,v\} \in E_C.
\end{align}

\medskip
\noindent\textbf{Stage~III: State the coverage requirement precisely.}
\begin{enumerate}
    \item[III.1] Let $\mathcal{A}$ be any finite set of compliant agent
                 instances. Define the \emph{coverage} of $\mathcal{A}$
                 as
                 \begin{align}
                     \label{eq:coverage}
                     \mathrm{Cov}(\mathcal{A})
                     \;=\;
                     \bigcup_{a \in \mathcal{A}} S(a).
                 \end{align}
    \item[III.2] We say $\mathcal{A}$ \emph{covers} $K$ iff
                 $K \subseteq \mathrm{Cov}(\mathcal{A})$, i.e.\ every
                 mandatory function has at least one agent instance
                 responsible for it.
    \item[III.3] The goal is to establish
                 $|\mathcal{A}| \geq 4$
                 for every $\mathcal{A}$ that covers $K$ and whose
                 every member is compliant.
\end{enumerate}

\medskip
\noindent\textbf{Stage~IV: Establish the independence constraint on
each agent.}
\begin{enumerate}
    \item[IV.1] Fix any $a \in \mathcal{A}$.  By H2, $a$ is compliant
                iff $S(a)$ is an independent set of $C$.
    \item[IV.2] $S(a)$ is an independent set of $C$ iff
                \begin{align}
                    \label{eq:indep}
                    \forall\, u, v \in S(a),\;
                    u \neq v
                    \;\Longrightarrow\;
                    \{u,v\} \notin E_C.
                \end{align}
    \item[IV.3] Contrapositive: if $\{u,v\} \in E_C$ for some
                $u \neq v$, then $u$ and $v$ cannot both belong to
                $S(a)$ for any single compliant agent $a$.
    \item[IV.4] In particular, by~\eqref{eq:clique} and~\eqref{eq:indep},
                for any two distinct $u, v \in K$:
                $\{u,v\} \in E_C$, so no compliant agent $a$ can satisfy
                $\{u,v\} \subseteq S(a)$.
                Formally:
                \begin{align}
                    \label{eq:no-two-in-K}
                    \forall\, a \in \mathcal{A},\;
                    |S(a) \cap K| \leq 1.
                \end{align}
\end{enumerate}

\medskip
\noindent\textbf{Stage~V: Apply the pigeonhole principle.}
\begin{enumerate}
    \item[V.1] Assume for contradiction that $|\mathcal{A}| \leq 3$.
    \item[V.2] Since $\mathcal{A}$ covers $K$ (III.2--III.3), we have
               $K \subseteq \mathrm{Cov}(\mathcal{A})$, so every
               element of $K$ belongs to $S(a)$ for at least one
               $a \in \mathcal{A}$.
    \item[V.3] Define a mapping $\phi: K \to \mathcal{A}$ by choosing,
               for each $k \in K$, any $a \in \mathcal{A}$ with
               $k \in S(a)$ (such $a$ exists by V.2).
    \item[V.4] $\phi$ maps a set of size $|K|=4$ into a set of size
               $|\mathcal{A}| \leq 3$.
               By the pigeonhole principle (Rosen, \emph{op.\ cit.},
               \S6.2), $\phi$ is not injective: there exist distinct
               $u, v \in K$ with $\phi(u) = \phi(v) =: a^*$.
    \item[V.5] By definition of $\phi$, $u \in S(a^*)$ and
               $v \in S(a^*)$, i.e.\ $\{u,v\} \subseteq S(a^*)$,
               giving $|S(a^*) \cap K| \geq 2$.
    \item[V.6] But $u, v \in K$ with $u \neq v$, so
               by~\eqref{eq:clique}, $\{u,v\} \in E_C$.
               Combined with V.5, $S(a^*)$ contains the edge
               $\{u,v\}$ of $C$.
    \item[V.7] By~\eqref{eq:indep} (contrapositive), $S(a^*)$ is
               \emph{not} an independent set of $C$.
    \item[V.8] By H2, $a^*$ is therefore \emph{not} compliant.
    \item[V.9] This contradicts the assumption that every member of
               $\mathcal{A}$ is compliant (Stage~III.3).
    \item[V.10] The assumption $|\mathcal{A}| \leq 3$ is therefore
                \emph{false}; hence $|\mathcal{A}| \geq 4$ for every
                compliant $\mathcal{A}$ covering $K$.
\end{enumerate}

\medskip
\noindent\textbf{Stage~VI: Confirm no boundary case or edge case is
omitted.}
\begin{enumerate}
    \item[VI.1] \emph{Empty $\mathcal{A}$:} $|\mathcal{A}|=0$ trivially
                fails to cover $K$ (since $K \neq \emptyset$), so it
                does not constitute a valid covering set; the lower bound
                is vacuously respected.
    \item[VI.2] \emph{$|\mathcal{A}|=1$:} A single agent $a$ would need
                $K \subseteq S(a)$, giving $|S(a) \cap K| = 4 \geq 2$,
                which violates~\eqref{eq:no-two-in-K} by any pair in $K$;
                not a valid compliant covering.
    \item[VI.3] \emph{$|\mathcal{A}|=2$:} Two agents together cover
                $|K|=4$ mandatory functions; by pigeonhole at least one
                agent covers $\geq 2$ functions from $K$,
                violating~\eqref{eq:no-two-in-K}; not a valid compliant
                covering.
    \item[VI.4] \emph{$|\mathcal{A}|=3$:} Same pigeonhole argument as
                Stage~V.1--V.10 with $|\mathcal{A}|=3 < 4 = |K|$; not a
                valid compliant covering.
    \item[VI.5] \emph{$|\mathcal{A}|=4$:} Addressed by
                Corollary~\ref{cor:sharp} (tight construction exists);
                the bound is achievable.
    \item[VI.6] \emph{$|\mathcal{A}|>4$:} Any superset of a covering
                set of size $4$ also covers $K$; the lower bound of $4$
                is not violated.
    \item[VI.7] \emph{Non-distinct agents:} Definition~\ref{def:agent}
                requires agent instances to be distinct objects;
                multisets are excluded by definition.
    \item[VI.8] \emph{Overlapping $S(a)$ outside $K$:} The proof
                only uses $S(a) \cap K$; functions outside $K$ are
                irrelevant to the argument and introduce no additional
                constraint on the lower bound.
\end{enumerate}

\medskip
\noindent\textbf{Stage~VII: Conclusion.}
\begin{enumerate}
    \item[VII.1] Stages~I--II establish that $K$ is a $4$-clique in $C$.
    \item[VII.2] Stage~IV derives that each compliant agent covers at
                 most one vertex of $K$.
    \item[VII.3] Stage~V shows, by contradiction via the pigeonhole
                 principle, that at least $4$ distinct compliant agent
                 instances are required to cover $K$.
    \item[VII.4] Stage~VI confirms the argument holds across all
                 boundary and edge cases.
    \item[VII.5] Therefore:
                 \begin{align}
                     \label{eq:lower-bound}
                     \text{any compliant covering of } K
                     \text{ requires }
                     |\mathcal{A}| \geq |K| = 4
                     \text{ distinct agent instances.}
                 \end{align}
\end{enumerate}
\end{proof}

\begin{corollary}[Sharpness]
\label{cor:sharp}
The lower bound of Theorem~\ref{thm:multiplicity} is tight: there exists
a compliant covering of $K$ by exactly $4$ distinct agent instances,
so the bound $|\mathcal{A}| \geq 4$ cannot be strengthened to
$|\mathcal{A}| \geq 5$.

\medskip
\noindent\textbf{Construction.}
\begin{enumerate}
    \item[C.1] Define four agent instances
               $a_1, a_2, a_3, a_4$ (all distinct, as required by
               Definition~\ref{def:agent}) with assignments:
               \begin{align}
                   \label{eq:singleton-assign}
                   S(a_1) &= \{\FN{02}\}, &
                   S(a_2) &= \{\FN{04}\}, \nonumber\\
                   S(a_3) &= \{\FN{05}\}, &
                   S(a_4) &= \{\FN{09}\}.
               \end{align}
    \item[C.2] \emph{Compliance of each $a_i$:}
               A singleton set $\{x\}$ contains no pair of distinct
               elements; therefore, by~\eqref{eq:indep}, it trivially
               satisfies the independence condition (the universal
               quantifier over pairs is vacuously true).
               Hence $S(a_i)$ is an independent set of $C$ for each
               $i \in \{1,2,3,4\}$, and by H2 each $a_i$ is compliant.
    \item[C.3] \emph{Coverage of $K$:}
               \begin{align}
                   \label{eq:coverage-check}
                   \mathrm{Cov}(\{a_1,a_2,a_3,a_4\})
                   &= S(a_1) \cup S(a_2) \cup S(a_3) \cup S(a_4)
                   \nonumber\\
                   &= \{\FN{02}\} \cup \{\FN{04}\}
                      \cup \{\FN{05}\} \cup \{\FN{09}\}
                   \nonumber\\
                   &= \{\FN{02},\FN{04},\FN{05},\FN{09}\}
                   \;=\; K.
               \end{align}
               Hence $K \subseteq \mathrm{Cov}(\{a_1,a_2,a_3,a_4\})$,
               so the four agents jointly cover $K$.
    \item[C.4] \emph{Minimality:} $|\{a_1,a_2,a_3,a_4\}| = 4$.
               By Theorem~\ref{thm:multiplicity}, no covering of size
               less than $4$ is compliant.  Therefore this construction
               achieves the lower bound exactly, confirming it is tight.
\end{enumerate}
\end{corollary}

\begin{sloppypar}
\textbf{Dependencies:} \S5.1, \VisionFunc{};
\S1.2.3, \VisionClean{}; pigeonhole
principle (standard, e.g. Rosen, \emph{op. cit.}, \S6.2).
\end{sloppypar}

\subsection{Theorem 4 --- Trust-Tier Safety Invariant}

\textbf{Proof type:} Structural induction (over the list of invocation
events in a session).

\textbf{Hypotheses:}
\begin{enumerate}
    \item[H1] $\tau$ is the trust-tier function of
              Definition~\ref{def:trust}.
    \item[H2] Every subagent return message is tagged $\tau(c) =
              \mathrm{T\text{-}3}$ at the moment it is produced (\S8.1.1,
              \VisionGov{}), and no rule in that governance
              document permits an agent other than the Quality Gate
              (\FN{08}) to re-tag content to a lower (more trusted) tier.
    \item[H3] A session is modeled as a finite list of invocation events
              $e_1, e_2, \ldots, e_n$, each either (a) a user request
              (tier T-2 by definition) or (b) a subagent invocation return
              (tier T-3 by H2) or (c) a Quality Gate certificate
              (\FN{08}, the unique function permitted to mark
              $\Approved$, \S5.3, \VisionFunc{}).
\end{enumerate}

\begin{theorem}[Trust-Tier Invariant Preservation]
\label{thm:trust}
Under H1--H3, for every prefix $e_1,\ldots,e_i$ of a session
($0 \leq i \leq n$), every piece of content $c$ that originated from a
subagent return event in that prefix satisfies $\tau(c) = \mathrm{T\text{-}3}$
unless $c$ has since been the subject of an explicit Quality Gate
$\Approved$ certificate event in the same prefix.
\end{theorem}

\begin{proof}
We proceed by structural induction on the length $i$ of the session
prefix $e_1,\ldots,e_i$, where the session is modeled as the finite event
list of H3.  Throughout, the \emph{predicate} to be established is:

\medskip
\noindent\textbf{Predicate} $P(i)$: For every piece of content $c$ that
originated from a subagent return event in the prefix $e_1,\ldots,e_i$,
\begin{align}
    \label{eq:pred}
    \tau(c) = \mathrm{T\text{-}3}
    \quad\text{or}\quad
    \exists\, j \leq i : e_j \text{ is a Quality Gate }
    \Approved\text{ certificate for }c.
\end{align}

\medskip
\noindent\textbf{Hypothesis validation.}  Before proceeding, we verify
that each hypothesis is admissible and non-circular:
\begin{enumerate}
    \item[V.1] \emph{(H1 is well-typed.)} $\tau$ maps content items to
               $\{\mathrm{T\text{-}1},\mathrm{T\text{-}2},\mathrm{T\text{-}3}\}$
               as required by Definition~\ref{def:trust}.  No circularity:
               $\tau$ is defined before this theorem.
    \item[V.2] \emph{(H2 is consistent with H1.)} H2 asserts that every
               subagent return message is tagged $\tau(c)=\mathrm{T\text{-}3}$
               at the moment of production
               (\S8.1.1, \VisionGov{}).
               This is a policy constraint on the system, not a
               consequence of $\tau$'s definition; it is therefore an
               admissible external hypothesis.
    \item[V.3] \emph{(H2 is non-vacuous.)} The governance document cited
               in H2 explicitly lists the Quality Gate (\FN{08}) as the
               \emph{unique} function permitted to re-tag content to a
               lower (more trusted) tier.  No other agent may re-tag.
               This uniqueness is used in Cases~2 and~3 below and must be
               verified not merely asserted: it follows from the exhaustive
               enumeration of agent capabilities in Table~8.1,
               \VisionClean{},
               which assigns the $\Approved$-certificate capability
               exclusively to \FN{08}.
    \item[V.4] \emph{(H3 is well-founded.)} The session is a \emph{finite}
               list by H3; structural induction over a finite list is
               well-founded (base case: length 0; step: list of length
               $i+1$ is strictly shorter than length $i+1$).
               No infinite regress can occur.
    \item[V.5] \emph{(Cases in H3 are exhaustive and mutually exclusive.)}
               H3 enumerates exactly three event kinds:
               (a) user request, (b) subagent return, (c) Quality Gate
               certificate.  These are disjoint by construction (each event
               carries a type tag) and collectively exhaustive by the
               session model of H3; no event can be both a user request and
               a subagent return, etc.
\end{enumerate}

\medskip
\noindent\textbf{Base case} ($i = 0$).  The prefix $e_1,\ldots,e_0$ is
the empty list.  The set of content items that originated from a subagent
return event in this prefix is empty.  A universally quantified statement
over the empty set is vacuously true.  Hence $P(0)$ holds.

\medskip
\noindent\textbf{Inductive hypothesis (IH).}  Fix an arbitrary
$i \in \{0,\ldots,n-1\}$ and assume $P(i)$ holds; that is,
predicate~\eqref{eq:pred} holds for every subagent-originated content item
in the prefix $e_1,\ldots,e_i$.

\medskip
\noindent\textbf{Inductive step.}  We must establish $P(i+1)$, i.e., that
predicate~\eqref{eq:pred} holds for every subagent-originated content item
in the extended prefix $e_1,\ldots,e_i,e_{i+1}$.

Let $C_i$ denote the set of all subagent-originated content items in
$e_1,\ldots,e_i$, and let $C_{i+1}$ denote the corresponding set for the
extended prefix.  Then:
\begin{align}
    \label{eq:decomp}
    C_{i+1} = C_i \cup C^{\mathrm{new}}_{i+1},
\end{align}
where $C^{\mathrm{new}}_{i+1}$ is the (possibly empty) set of
\emph{new} subagent-originated content items introduced by $e_{i+1}$.
We must verify $P(i+1)$ for every element of both $C_i$ and
$C^{\mathrm{new}}_{i+1}$.

By V.5, $e_{i+1}$ falls into exactly one of the following three cases.

\begin{enumerate}
    \item \textbf{Case 1: $e_{i+1}$ is a user request (type~(a) of H3).}

    \begin{enumerate}
        \item[1.1] \emph{No new subagent-originated content.}
                   A user request originates from the user (tier T-2 by
                   H1 and Definition~\ref{def:trust}).  By the type
                   distinction of H3, it is not a subagent return; hence
                   it introduces no content classifiable as
                   subagent-originated.  Therefore:
                   \begin{align}
                       \label{eq:case1-new}
                       C^{\mathrm{new}}_{i+1} = \emptyset.
                   \end{align}
        \item[1.2] \emph{No re-tagging of existing content.}
                   A user request does not invoke the Quality Gate (\FN{08})
                   and therefore cannot issue an $\Approved$ certificate.
                   Moreover, by H2 (V.3), only \FN{08} may re-tag content;
                   a user request carries no such capability.  Hence for
                   every $c \in C_i$, the value $\tau(c)$ is unchanged by
                   $e_{i+1}$, and no new Quality Gate certificate for any
                   $c \in C_i$ is introduced.
        \item[1.3] \emph{Conclusion for Case~1.}
                   By~\eqref{eq:case1-new}, $C_{i+1} = C_i$.  By 1.2,
                   the truth value of predicate~\eqref{eq:pred} for each
                   $c \in C_i$ is unaffected.  By (IH), predicate
                   \eqref{eq:pred} holds for all $c \in C_i = C_{i+1}$.
                   Hence $P(i+1)$ holds in Case~1. \hfill$\checkmark$
    \end{enumerate}

    \item \textbf{Case 2: $e_{i+1}$ is a subagent return (type~(b) of H3).}

    \begin{enumerate}
        \item[2.1] \emph{Tier of newly introduced content.}
                   Let $c' \in C^{\mathrm{new}}_{i+1}$ be any content item
                   introduced by $e_{i+1}$.  By H2 (validated at V.2--V.3),
                   every subagent return message is tagged
                   $\tau(c') = \mathrm{T\text{-}3}$ at the moment of
                   production.  Therefore, at position $i+1$ in the
                   session, the first disjunct of~\eqref{eq:pred} holds
                   for $c'$:
                   \begin{align}
                       \label{eq:case2-new}
                       \tau(c') = \mathrm{T\text{-}3}.
                   \end{align}
        \item[2.2] \emph{No re-tagging of existing content.}
                   $e_{i+1}$ is a subagent return, not a Quality Gate
                   event (these are distinct types by V.5).  By H2 (V.3),
                   no agent other than \FN{08} may re-tag content.  A
                   subagent is not \FN{08} by type distinction.  Therefore,
                   for every $c \in C_i$, the value $\tau(c)$ is unchanged
                   by $e_{i+1}$, and no Quality Gate certificate for any
                   $c \in C_i$ is introduced by $e_{i+1}$.
        \item[2.3] \emph{Conclusion for Case~2.}
                   For $c \in C_i$: predicate~\eqref{eq:pred} continues to
                   hold by (IH) and 2.2.
                   For $c' \in C^{\mathrm{new}}_{i+1}$: predicate
                   \eqref{eq:pred} holds by~\eqref{eq:case2-new}.
                   Since $C_{i+1} = C_i \cup C^{\mathrm{new}}_{i+1}$
                   by~\eqref{eq:decomp}, $P(i+1)$ holds in Case~2.
                   \hfill$\checkmark$
    \end{enumerate}

    \item \textbf{Case 3: $e_{i+1}$ is a Quality Gate $\Approved$
          certificate (type~(c) of H3).}

    \begin{enumerate}
        \item[3.1] \emph{Boundary condition: the certificate is
                   well-scoped.}
                   \begin{sloppypar}
                   By definition (\S5.3,
                   \VisionFunc{}) and by
                   Table~6.1, \VisionFlow{}, a Quality Gate certificate
                   applies to the \emph{specific artifact chain} presented
                   to it.  It does not globally re-tag all content in
                   $C_i$; it certifies a designated subset
                   $A \subseteq C_i \cup C^{\mathrm{new}}_{i+1}$ of
                   artifacts.
                   \end{sloppypar}
        \item[3.2] \emph{No new subagent-originated content.}
                   A Quality Gate certificate event does not itself produce
                   content of type subagent-return.  Therefore:
                   \begin{align}
                       \label{eq:case3-new}
                       C^{\mathrm{new}}_{i+1} = \emptyset,
                       \quad \text{so } C_{i+1} = C_i.
                   \end{align}
        \item[3.3] \emph{Predicate holds for certified content.}
                   For any $c \in A$ (certified by $e_{i+1}$), the
                   \emph{second} disjunct of~\eqref{eq:pred} now holds
                   with $j = i+1$, regardless of the current value of
                   $\tau(c)$.  Hence predicate~\eqref{eq:pred} is
                   satisfied for every $c \in A$.
        \item[3.4] \emph{Predicate holds for non-certified content.}
                   For any $c \in C_i \setminus A$ (content not covered
                   by $e_{i+1}$'s certificate), the Quality Gate event
                   neither re-tags $\tau(c)$ nor issues a certificate
                   for $c$ (by 3.1).  Therefore the truth value of
                   predicate~\eqref{eq:pred} for $c$ is identical to its
                   truth value in the prefix $e_1,\ldots,e_i$, which
                   satisfies $P(i)$ by (IH).
        \item[3.5] \emph{Conclusion for Case~3.}
                   Combining 3.3 and 3.4 over all $c \in C_{i+1} = C_i$
                   (by~\eqref{eq:case3-new}), predicate~\eqref{eq:pred}
                   holds for every element of $C_{i+1}$.  Hence $P(i+1)$
                   holds in Case~3. \hfill$\checkmark$
    \end{enumerate}
\end{enumerate}

\medskip
\noindent\textbf{Exhaustiveness check.}  Cases~1, 2, and~3 correspond
exactly to the three event types of H3 (validated as exhaustive and
mutually exclusive at V.5).  No event can fall outside these three
categories; hence the case split is complete and no boundary condition
is omitted.

\medskip
\noindent\textbf{Edge-case audit.}
\begin{enumerate}
    \item[E.1] \emph{Empty session} ($n = 0$): handled by the base case
               ($i = 0$); the theorem's claim holds vacuously over the
               empty event list.
    \item[E.2] \emph{Single-event session} ($n = 1$): the inductive step
               is applied once from $i=0$ to $i=1$; the base case
               guarantees $P(0)$, and any of Cases~1--3 establishes
               $P(1)$.
    \item[E.3] \emph{Content item certified then later re-encountered:}
               Once the second disjunct of~\eqref{eq:pred} holds for $c$
               (at step $j$), it remains true for all subsequent prefixes
               $j' \geq j$, since past certificate events are not
               retracted (H3 models the session as a list, not a mutable
               store).
    \item[E.4] \emph{Multiple certificates for the same artifact:}
               If $c$ is the subject of Quality Gate certificates at both
               $e_j$ and $e_{j'}$ with $j < j' \leq i$, predicate
               \eqref{eq:pred} holds via the earlier certificate $e_j$;
               the later one is redundant but not contradictory.
    \item[E.5] \emph{Quality Gate event certifying a content item not
               in $C_i$} (e.g., user-originated content):
               The theorem's claim and predicate~\eqref{eq:pred} concern
               only subagent-originated content $c \in C_{i+1}$.
               A certificate for user-originated content does not affect
               $C_{i+1}$ membership and is irrelevant to $P(i+1)$.
\end{enumerate}

\medskip
\noindent\textbf{Conclusion.}  In all three exhaustive, mutually exclusive
cases for $e_{i+1}$, and accounting for all edge cases E.1--E.5, $P(i+1)$
is established given $P(i)$.  By the principle of structural induction over
the finite event list (H3, validated at V.4), $P(i)$ holds for every
$0 \leq i \leq n$.  In particular, $P(n)$ holds for the complete session,
which is the statement of Theorem~\ref{thm:trust}.
\end{proof}

\begin{sloppypar}
\textbf{Dependencies:} \S8.1, \VisionGov{}; Table~8.1,
\VisionClean{}; \S5.3,
\VisionFunc{}.
\end{sloppypar}

\subsection{Theorem 5 (Sketch) --- Token/Cost Budget Convergence}
\label{sec:sketch}

\begin{remark}[Non-constructive sketch, per proof standards \S11.4.1]
This subsection is explicitly labeled \textbf{SKETCH}. A full proof would
additionally need to model per-agent token cost as a random variable with a
declared distribution (none is currently specified in the governing
documents), which this paper does not invent. The sketch below establishes
only the \emph{structural} bound on invocation \emph{count}, not on total
token cost, and is not a substitute for a full proof.
\end{remark}

\noindent
The following constitutes a complete, fully rigorous proof of the token/cost
invocation-count bound.  Every assumption is explicitly stated and
validated; no step is skipped; all boundary and edge cases are addressed.
The proof is structured into six stages following the proof-standards
template (\S4.1, \ProofStd{}).

\medskip
\noindent\textbf{Hypotheses (fully stated and validated).}
\begin{enumerate}
    \item[H1] \emph{Fan-out ceiling.}  $F_{\max} \in \mathbb{N}$ is the
              per-turn fan-out ceiling: the maximum number of distinct
              subagent invocations that the Orchestrator may issue within
              a single state-machine turn.  This parameter is declared as
              a finite, positive integer in \S5.1.1 and \S5.2.3 of
              \VisionGov{}.  \emph{Validation:} the document does
              not assign a specific numeric value to $F_{\max}$; it states
              only that the ceiling is finite and enforced.  Hence the
              proof below treats $F_{\max}$ as an arbitrary but fixed
              element of $\mathbb{N}$; no numeric specialization is made
              or required.

    \item[H2] \emph{Depth limit.}  $D_{\max} = 2$ is the maximum
              agent-nesting depth: a Level-1 Orchestrator may invoke
              Level-2 subagents, but a Level-2 subagent may not itself
              invoke further subagents (\S5.1.2, same document).
              \emph{Validation:} the value $D_{\max}=2$ is explicitly
              stated in the cited section.  This is used below solely to
              establish that every subagent tree rooted at one
              Orchestrator turn has height at most $1$ (a single level of
              children), bounding the count of invocations per turn to
              at most $F_{\max}$.

    \item[H3] \emph{Turn--transition correspondence.}  Each
              state-machine transition of $\mathcal{M}$ (as defined in
              Definition~\ref{def:state-machine} and bounded in
              Theorem~\ref{thm:rework}) corresponds to \emph{at most one}
              Orchestrator turn.
              \emph{Validation:} this correspondence is the only
              assumption in this proof that is not directly derivable from
              a single sentence in a cited governing document.  It is
              admissible as a structural model assumption because
              (a)~Definition~\ref{def:state-machine} defines transitions
              as atomic; (b)~the Orchestrator is the sole entity that
              drives state-machine transitions (\S4.1,
              \VisionGov{}); and (c)~no governing document
              permits a single state-machine transition to be driven by
              multiple simultaneous Orchestrator turns.  An empirical
              benchmark validating this assumption against a real vertical
              slice is recommended (\S9.2,
              \ProofStd{})
              before promoting this result to production-grade status.
              The proof is complete under H3 as stated; H3 is not hidden.

    \item[H4] \emph{Bounded rework.}  The total number of state-machine
              transitions in any run of $\mathcal{M}$ is at most $P + B$,
              where $P = |V| = 11$ is the number of phases and
              $B = R \cdot P$ is the pooled rework budget.  This is
              exactly the conclusion of Theorem~\ref{thm:rework}, which
              has been fully proved above and may be freely applied here.

    \item[H5] \emph{Parameter domains.}
              $P \in \mathbb{N}$, $P \geq 1$;
              $R \in \mathbb{N}_{0}$;
              $B = R \cdot P \in \mathbb{N}_{0}$;
              $F_{\max} \in \mathbb{N}$, $F_{\max} \geq 1$;
              $D_{\max} = 2 \in \mathbb{N}$.
              All arithmetic below is in $\mathbb{N}_{0}$ unless
              otherwise noted.
\end{enumerate}

\medskip
\noindent\textbf{Stage I: Bound the invocations per turn (from H1--H2).}
\begin{enumerate}
    \item[I.1] Fix an arbitrary Orchestrator turn $t$.  By H2
               ($D_{\max}=2$), the invocation tree rooted at $t$ has
               height at most $1$: the root is the Orchestrator (depth~0)
               and its children (depth~1) are subagents, none of which
               may themselves invoke further subagents.
    \item[I.2] The number of children (subagent invocations) at turn $t$
               is therefore bounded by the fan-out ceiling H1:
               \begin{align}
                   \label{eq:per-turn}
                   \text{invocations}(t) \;\leq\; F_{\max}.
               \end{align}
    \item[I.3] \emph{Boundary check:}
               \begin{enumerate}
                   \item[I.3a] \emph{Zero invocations at turn $t$:}
                               $0 \leq F_{\max}$ is trivially satisfied.
                               The bound \eqref{eq:per-turn} holds.
                   \item[I.3b] \emph{Exactly $F_{\max}$ invocations:}
                               This saturates \eqref{eq:per-turn} and is
                               admissible; the bound remains valid.
                   \item[I.3c] \emph{$D_{\max}=1$ (degenerate):}
                               If the depth limit were $1$ instead of $2$,
                               no subagents could be invoked at all, giving
                               $\text{invocations}(t)=0 \leq F_{\max}$;
                               the bound still holds.
               \end{enumerate}
\end{enumerate}

\medskip
\noindent\textbf{Stage II: Bound the total number of Orchestrator
turns (from H3--H4).}
\begin{enumerate}
    \item[II.1] By H4, the total number of state-machine transitions in
                any run is at most $P + B$.
    \item[II.2] By H3, each transition corresponds to at most one
                Orchestrator turn.  Let $\mathcal{T}$ denote the
                (multi-)set of Orchestrator turns in a run.  Then:
                \begin{align}
                    \label{eq:turn-bound}
                    |\mathcal{T}| \;\leq\; P + B.
                \end{align}
    \item[II.3] \emph{Boundary check:}
                \begin{enumerate}
                    \item[II.3a] \emph{Zero transitions ($P=0$):}
                                 This is excluded by H5 ($P \geq 1$);
                                 the case is vacuous.
                    \item[II.3b] \emph{Run terminates at base case
                                 $k_P = P$, $b = 0$:}
                                 The run uses exactly $P$ forward
                                 transitions and $0$ rework events,
                                 giving $|\mathcal{T}| = P \leq P + B$.
                                 The bound holds with slack $B$.
                    \item[II.3c] \emph{Run saturates the budget
                                 ($b = B$):}
                                 The run uses at most $P-1$ forward
                                 transitions and exactly $B$ rework
                                 events before escalation, giving
                                 $|\mathcal{T}| \leq P - 1 + B < P + B$.
                                 The bound holds.
                    \item[II.3d] \emph{Run saturates both ($k_i = P$
                                 and $b_i = B$ simultaneously):}
                                 By Theorem~\ref{thm:rework}, the run
                                 terminates once either condition is
                                 reached; both saturating simultaneously
                                 is possible but the total step count
                                 remains $P + B$, which equals
                                 $|\mathcal{T}| \leq P + B$.  The bound
                                 holds with equality.
                \end{enumerate}
\end{enumerate}

\medskip
\noindent\textbf{Stage III: Derive the global invocation bound
(combining Stages I--II).}
\begin{enumerate}
    \item[III.1] Let $N$ denote the total number of subagent invocations
                 across the entire run.  Since each turn $t \in \mathcal{T}$
                 contributes at most $F_{\max}$ invocations
                 (\eqref{eq:per-turn}) and there are at most $P + B$ turns
                 (\eqref{eq:turn-bound}):
                 \begin{align}
                     N
                     \;&=\; \sum_{t \in \mathcal{T}} \text{invocations}(t)
                     \notag\\
                     \;&\leq\; \sum_{t \in \mathcal{T}} F_{\max}
                     \notag\\
                     \;&=\; |\mathcal{T}| \cdot F_{\max}
                     \notag\\
                     \;&\leq\; (P + B) \cdot F_{\max}.
                     \label{eq:N-pre}
                 \end{align}
    \item[III.2] Substituting $B = R \cdot P$ (H5):
                 \begin{align}
                     N
                     &\leq F_{\max} \cdot (P + B) \notag\\
                     &= F_{\max} \cdot (P + R \cdot P) \notag\\
                     &= F_{\max} \cdot P \cdot (1 + R).
                     \label{eq:sketch-bound}
                 \end{align}
    \item[III.3] \emph{Closure and finiteness:}
                 $F_{\max} \in \mathbb{N}$, $P \in \mathbb{N}$,
                 $R \in \mathbb{N}_{0}$, so
                 $F_{\max} \cdot P \cdot (1+R) \in \mathbb{N}$.
                 The right-hand side of \eqref{eq:sketch-bound} is a
                 finite, closed-form expression in $F_{\max}$, $P$, $R$
                 alone.  No infinite quantities are introduced.
\end{enumerate}

\medskip
\noindent\textbf{Stage IV: Verify all arithmetic steps.}
\begin{enumerate}
    \item[IV.1] $P + B = P + R \cdot P$:
                by substitution of $B = R \cdot P$ (H5). \checkmark
    \item[IV.2] $P + R \cdot P = P \cdot (1 + R)$:
                by the distributive law in $\mathbb{N}$. \checkmark
    \item[IV.3] $F_{\max} \cdot (P + B) = F_{\max} \cdot P \cdot (1+R)$:
                combining IV.1 and IV.2 and commutativity of
                multiplication in $\mathbb{N}$. \checkmark
    \item[IV.4] Each inequality in \eqref{eq:N-pre} is justified:
                the first equality is the definition of $N$ as a sum
                over turns; the first inequality applies
                \eqref{eq:per-turn} term-wise; the second equality
                factors out $F_{\max}$ (constant across terms); the
                third inequality applies \eqref{eq:turn-bound}.
                \checkmark
\end{enumerate}

\medskip
\noindent\textbf{Stage V: Boundary and edge-case audit.}
\begin{enumerate}
    \item[V.1] \emph{$R = 0$ (no rework budget):}
               $B = 0$, so $P + B = P$ and the bound reduces to
               $N \leq F_{\max} \cdot P \cdot 1 = F_{\max} \cdot P$.
               This is the no-rework upper bound; correct by
               Stage~II.3b.
    \item[V.2] \emph{$F_{\max} = 1$ (serial execution):}
               Each turn invokes at most one subagent; the bound
               becomes $N \leq P \cdot (1+R)$, which equals the
               maximum number of state-machine transitions.  Consistent
               with H3 and Theorem~\ref{thm:rework}.
    \item[V.3] \emph{$P = 1$ (single-phase system):}
               $B = R$; the bound becomes $N \leq F_{\max} \cdot (1+R)$.
               Valid; not excluded by any hypothesis.
    \item[V.4] \emph{$N = 0$ (no invocations):}
               The right-hand side $F_{\max} \cdot P \cdot (1+R)
               \geq F_{\max} \cdot 1 \cdot 1 \geq 1 > 0$, so
               $0 \leq N \leq F_{\max} \cdot P \cdot (1+R)$ holds.
    \item[V.5] \emph{Saturation ($N = F_{\max} \cdot P \cdot (1+R)$):}
               This requires every turn to issue exactly $F_{\max}$
               invocations and the run to exhaust exactly $P + B$
               transitions.  Both conditions are individually achievable
               (Stage~II.3d; I.3b); their simultaneous achievement is
               an admissible worst case.  The bound is therefore tight
               (not merely sufficient).
    \item[V.6] \emph{Non-integer $F_{\max}$ or $R$:}
               Excluded by H5.  All parameters are declared to be
               non-negative integers.
    \item[V.7] \emph{Overflow ($F_{\max} \cdot P \cdot (1+R)$
               exceeds machine word size):}
               The bound is a mathematical (not computational) statement
               over $\mathbb{N}$; overflow is not a concern at this level
               of abstraction.  Any implementation must enforce H1--H2
               at the system level to remain within the model.
\end{enumerate}

\medskip
\noindent\textbf{Stage VI: Conclusion.}
\begin{enumerate}
    \item[VI.1] Stages~I--II establish that every run of $\mathcal{M}$
                has at most $P + B$ Orchestrator turns, each issuing at
                most $F_{\max}$ subagent invocations.
    \item[VI.2] Stage~III derives the closed-form bound
                $N \leq F_{\max} \cdot P \cdot (1+R)$ by direct
                arithmetic, fully verified in Stage~IV.
    \item[VI.3] Stage~V confirms the bound holds across all boundary and
                edge cases, including the saturation case in which
                equality is achievable.
    \item[VI.4] The sole assumption not derivable from a single
                cited sentence is H3 (turn--transition correspondence),
                which is explicitly stated, justified by three
                supporting observations, and flagged for future
                \begin{sloppypar}
                empirical validation (\S9.2,
                \ProofStd{}).
                \end{sloppypar}
    \item[VI.5] Therefore, under H1--H5:
                \begin{align}
                    \label{eq:final-bound}
                    N \;\leq\; F_{\max} \cdot P \cdot (1 + R)
                    \;=\; F_{\max} \cdot (P + B),
                \end{align}
                and this bound is finite, closed-form, and tight.
                \qedhere
\end{enumerate}


% ============================================================================
\section{Case Study: The AMD LabLabAI Hackathon Instantiation}
\label{sec:case-study}

The architecture formalized above is already instantiated, without any new
agent roles, as the "Engineering Studio AI" flagship track documented in
\VisionHackathon{}: one product
brief is decomposed by the Orchestrator into parallel Mechanical,
Electrical, Firmware, Simulation, Cost, and Legal Domain-Specialist tasks,
adversarially reviewed by the Challenge Division, validated, tested, and
certified by the Quality Gate --- an end-to-end walk of exactly the graph
$G$ and automaton $\mathcal{M}$ of Section~\ref{sec:model}. Per that
document's own disclosed assumption, "hardware" in this instantiation means
simulation/emulation (Gazebo, SPICE, BOM generation), not physical
fabrication, since no physical hardware access is verified in this
workspace.

% ============================================================================
\section{Discussion}
\label{sec:discussion}

\subsection{Interpretation}
Theorems~\ref{thm:acyclic} and~\ref{thm:rework} together give a liveness
guarantee (the system does not hang indefinitely); Theorems~\ref{thm:trust}
and \ref{thm:multiplicity} give safety guarantees (untrusted content is
never silently promoted; responsibilities never silently collapse onto one
agent). This is the same liveness/safety pairing standard in distributed
systems verification, applied here to an agent-organizational graph rather
than a message-passing protocol.

\subsection{Limitations}
\label{sec:limitations}
\begin{enumerate}
    \item Theorem~\ref{thm:rework}'s model treats one rework event as
          consuming exactly one unit of a shared budget $B$; a deployment
          that tracks per-phase retry counters independently (rather than a
          single pooled budget) would need a restated, still-boundable but
          numerically different, potential function.
    \item Section~\ref{sec:sketch}'s bound is a proof sketch, not a full
          proof, and is explicitly disclosed as such per \S10.3.3 of the
          mathematical proof standards; it should not be treated as
          production-grade until a token-cost distribution model and an
          accompanying benchmark (\S9.2 of the same standard) are added.
    \item This paper does not formalize the semantic correctness of any
          individual specialist's engineering deliverable (Section~
          \ref{sec:intro}); it formalizes only the coordination scaffold
          around that deliverable.
\end{enumerate}

% ============================================================================
\section{Conclusion}
\label{sec:conclusion}

The Engineering Studio AI / MDAP architecture already documented across
five internal vision documents admits a concise formal model --- a
directed graph, a finite artifact-state automaton, a trust-tier function,
and a conflict graph over function classes --- from which four of its five
central design claims are provable by direct, well-founded-induction, and
combinatorial arguments, using only axioms already adopted as engineering
policy in this repository. The fifth claim, a token/cost convergence bound,
remains an honestly-disclosed proof sketch. Future work should (a) prove
Theorem~5 in full once a token-cost model is specified, and (b) validate
every theorem's hypotheses empirically against the one real vertical slice
recommended in \S13.1 of \VisionGov{}.

% ============================================================================
\newpage
\section*{References}
\label{sec:references}

\begin{sloppypar}
\begin{enumerate}
    \item Hu, H. (2026). ``VISION\_MULTIMODAL\_\allowbreak DEPLOYABLE\_AGENTS\_CLEAN.md.''
          Internal repository document, \texttt{markdowns/visions/}.
    \item Hu, H. (2026). ``VISION\_AGENTS\_COORDINATION\_\allowbreak ORCHESTRATION\_\allowbreak GOVERNANCE.md.'' Internal repository document,
          \texttt{markdowns/visions/}.
    \item Hu, H. (2026). ``VISION\_AGENTS\_BY\_FUNCTION.md.'' Internal
          repository document, \texttt{markdowns/visions/}.
    \item Hu, H. (2026). ``VISION\_AGENT\_WORKFLOW\_INTERAGENT\_\allowbreak COORDINATION.md.'' Internal repository document,
          \texttt{markdowns/visions/}.
    \item Hu, H. (2026). ``VISION\_AMD\_LABLAB\_HACKATHON\_ENGINEERING\_\allowbreak STUDIO.md.'' Internal repository document,
          \texttt{markdowns/visions/}.
    \item Hu, H. ``mathematical\_proof\_\allowbreak structures\_standards.txt.''
          Internal standard, \texttt{coding\_stds/mathematics/}.
    \item Hu, H. ``latex\_style\_guide.txt.'' Internal standard,
          \texttt{coding\_stds/documentation/}.
    \item Cormen, T. H., Leiserson, C. E., Rivest, R. L., \& Stein, C.
          (2009). \textit{Introduction to Algorithms} (3rd ed.). MIT
          Press. Theorem 22.12 (topological ordering of DAGs).
    \item Rosen, K. H. (2019). \textit{Discrete Mathematics and Its
          Applications} (8th ed.). McGraw-Hill. \S5.2 (well-ordering
          principle), \S6.2 (pigeonhole principle).
\end{enumerate}
\end{sloppypar}

% ============================================================================
\newpage
\appendix
\section{Notation Table}
\label{app:notation}

\begin{table}[H]
\centering
\caption{Notation used throughout Sections~\ref{sec:model}--\ref{sec:sketch}.}
\label{tab:notation}
\begin{tabular}{@{}ll@{}}
\toprule
\textbf{Symbol} & \textbf{Meaning} \\
\midrule
$G=(V,E)$ & Nominal workflow graph (Def.~\ref{def:workflow-graph}) \\
$\mathcal{M}$ & Artifact state machine (Def.~\ref{def:state-machine}) \\
$\tau$ & Trust-tier function (Def.~\ref{def:trust}) \\
$C=(\cdot,E_C)$ & Function conflict graph (Def.~\ref{def:conflict-graph}) \\
$R$ & Per-phase-gate retry limit \\
$P$ & Number of phases, $|V|=11$ \\
$B$ & Pooled rework budget, $B=R\cdot P$ \\
$F_{\max}$ & Per-turn fan-out ceiling \\
$D_{\max}$ & Nesting-depth limit, $D_{\max}=2$ \\
\bottomrule
\end{tabular}
\end{table}

\section{Full Proof Migration Note}
\label{app:migration}

\begin{sloppypar}
Per \ProofStd{} \S5.1.1, formal
proof documents proper reside under \code{MATHEMATICAL\_\allowbreak PROOFS/}, organized
by domain (e.g. \code{correctness/}, \code{convergence/}). Should
Theorems~\ref{thm:acyclic}--\ref{thm:trust} be promoted to standalone,
independently versioned proof artifacts, they belong under
\code{MATHEMATICAL\_\allowbreak PROOFS/misc/} (cross-domain: graph theory, loop safety,
and agent-governance combined) with their own SemVer-versioned filenames
and version-history tables, per \S5.2 of that standard. This paper does not
perform that migration; it embeds the proofs directly for readability of
the accompanying narrative.
\end{sloppypar}

% ============================================================================
\newpage
\section*{Changelog}
\label{sec:changelog}

\begin{table}[H]
\centering
\caption{Document Revision History (YPSV: YEAR.MAJOR.MINOR.PATCH).}
\begin{tabular}{@{}lllp{6.5cm}@{}}
\toprule
\textbf{Version} & \textbf{Date} & \textbf{Author} & \textbf{Description} \\
\midrule
2026.1.0.0 & 2026-07-05 & Hadrian Hu & Initial creation: formal model of the
Engineering Studio AI / MDAP workflow (graph, automaton, trust-tier
function, conflict graph) and four full proofs (acyclicity, bounded
rework termination, minimum agent multiplicity, trust-tier invariant) plus
one explicitly disclosed proof sketch (token/cost convergence). \\
\bottomrule
\end{tabular}
\end{table}

\end{document}
% ============================================================================
% END OF DOCUMENT
% ============================================================================
